\section*{Solutions to Homework 10}

\begin{lstlisting}
$ bash codeHints/L10/setup_hw10.sh
$ cd practice/hw10
$ git reflog | grep "Important work"            # 10.1
0b30b44 HEAD@{18}: commit: Important work
$ git branch rescued 0b30b44                      # use the ID you see
$ git switch messy                                # 10.2
$ git reset --soft main
$ git commit -m "Add feature"
$ git switch history                              # 10.3
$ git bisect start
$ git bisect bad
$ git bisect good $(git rev-list --max-parents=0 history)
$ git bisect run bash test.sh
... is the first bad commit
$ git bisect reset
$ git switch main                                 # 10.4
$ echo "work in progress" >> notes.txt
$ git stash
$ git switch urgent
$ sed -i.bak 's/teh/the/' app.txt && rm app.txt.bak     # or edit by hand
$ git commit -am "Fix urgent bug"
$ git switch main
$ git stash pop
$ git status --short
 M notes.txt
\end{lstlisting}

For 10.5 (the three IDs are examples, yours differ except the hash of \code{hello}):

\begin{lstlisting}
$ git cat-file -p main | head -1
tree 8cfff4b...
$ echo hello | git hash-object --stdin
ce013625030ba8dba906f756967f9e9ca394464a
$ cat .git/refs/heads/main
1689350...
$ printf 'bad-commit=e238140\ntree=8cfff4b\nhello=ce01362\nref-file=1689350\n' > ../hw10_answers.txt
\end{lstlisting}

\textbf{10.1.} The reset moved \code{main} back, but the reflog remembers every position of \code{HEAD}. Your reflog is long because it also contains the history of the setup; \code{grep} finds the line. \code{git branch rescued <id>} only adds a label, so \code{main} stays where it was. (If you used \code{git reset --hard <id>} on \code{main}, the rescue works but \code{main} moved, and the checker complains.)

\textbf{10.2.} \code{git reset --soft main} moves \code{messy} back to \code{main}'s commit but keeps all the changes \emph{staged}, so a single \code{git commit} makes one commit with exactly the same files as the three wip commits (the checker compares the file contents).

\textbf{10.3.} \code{bisect run} tests about $\log_2 12 \approx 4$ commits. The first bad commit is \code{Step 7}. Do not forget \code{git bisect reset}: otherwise the repository stays in bisect mode and the checker reports it.

\textbf{10.4.} A stash lets you switch branches with a clean working directory. After the pop, your uncommitted line is back in \code{notes.txt}, and the commit \code{Fix urgent bug} lives on \code{urgent}, not on \code{main}. If \code{git stash pop} reports a conflict, you changed the same lines on both sides; resolve it as in Lecture 5.

\textbf{10.5.} \code{git cat-file -p main} prints a commit: its first line gives the ID of its tree. The hash of \code{hello} is the same on every computer (\code{ce01362\dots}). The file \code{.git/refs/heads/main} holds the full ID of \code{main}'s commit, so its first 7 characters equal \code{git rev-parse --short=7 main}. (If the file does not exist, Git has packed the reference into \code{.git/packed-refs}; use \code{git rev-parse main} instead.)

\Needspace{8\baselineskip}
\textbf{Common mistakes.}
\begin{itemize}
  \item Moving \code{main} to the lost commit in 10.1 instead of creating a branch.
  \item Forgetting \code{git bisect reset}; or running bisect while having uncommitted changes (stash or commit first).
  \item Committing the \code{work in progress} line instead of stashing it in 10.4: then \code{main} has a new commit and the check ``only \code{ M notes.txt}'' fails.
  \item Writing the \emph{whole} 40-character ID is fine, but the checker also needs at least 7 characters.
\end{itemize}
